\documentclass[14pt, a4paper, landscape]{extarticle}

\usepackage{xcolor}
\usepackage{geometry}
\usepackage{graphicx}
\usepackage{booktabs}
\usepackage{array}
\usepackage{longtable}
\usepackage{titlesec}
\usepackage{enumitem}
\usepackage{tcolorbox}
\usepackage{fancyhdr}
\usepackage{tikz}

\geometry{margin=1cm}

\usepackage{xepersian}
\settextfont{Vazirmatn}
\setlatintextfont{Times New Roman}
\setdigitfont{Vazirmatn}

\definecolor{primary}{RGB}{139, 90, 43}
\definecolor{accent}{RGB}{224, 122, 63}
\definecolor{bg}{RGB}{245, 240, 232}
\definecolor{darkbg}{RGB}{92, 58, 26}

\pagestyle{fancy}
\fancyhf{}
\fancyhead[R]{\small\textcolor{primary}{مهندسی نرم‌افزار — جلسهٔ ۱}}
\fancyfoot[C]{\textcolor{primary}{\thepage}}
\renewcommand{\headrulewidth}{0.5pt}
\setlength{\headheight}{15pt}

\setlength{\parindent}{0pt}
\setlength{\parskip}{0.5em}

\newcommand{\slide}[1]{%
    \newpage
    \begin{tcolorbox}[
        colback=bg,
        colframe=primary,
        boxrule=1.5pt,
        arc=4mm,
        left=8mm, right=8mm, top=6mm, bottom=6mm,
        width=\textwidth,
        height=0.85\textheight
    ]
    #1
    \end{tcolorbox}
}

\newcommand{\slidetitle}[1]{%
    \begin{center}
        {\LARGE\bfseries\textcolor{primary}{#1}}
    \end{center}
    \vspace{0.5cm}
}

\newcommand{\myaccent}[1]{\textcolor{accent}{#1}}

\begin{document}

\pagestyle{empty}

\begin{tikzpicture}[remember picture, overlay]
    \fill[darkbg] (current page.north west) rectangle (current page.south east);
\end{tikzpicture}

\begin{center}
    \vspace*{4cm}
    {\color{white}\Huge\bfseries مهندسی نرم‌افزار}
    
    \vspace{1cm}
    
    {\color{accent}\LARGE جلسهٔ ۱ — معرفی درس}
    
    \vspace{2cm}
    
    {\color{white}\Large ابراهیم فضلی}
    
    \vspace{0.5cm}
    
    {\color{white}\large ترم ۱۴۰۴}
\end{center}

\pagestyle{fancy}

\slide{
    \slidetitle{یک تصویر، یک سؤال}
    \vspace{1cm}
    \begin{center}
        {\Large\bfseries نرم‌افزار هم می‌تواند شکست بخورد.}
        \vspace{1.5cm}
        {\Huge\bfseries\myaccent{چرا؟}}
    \end{center}
}

\slide{
    \slidetitle{آمار تکان‌دهنده}
    \vspace{1cm}
    \begin{center}
        {\fontsize{80}{90}\selectfont\bfseries\myaccent{۷۰٪}}
        \vspace{1cm}
        {\Large پروژه‌های نرم‌افزاری شکست می‌خورند،}
        {\Large دیر تحویل می‌شوند یا از بودجه فراتر می‌روند.}
        \vspace{1.5cm}
        {\small منبع: Standish Group Chaos Report}
    \end{center}
}

\slide{
    \slidetitle{کدنویس vs مهندس نرم‌افزار}
    \vspace{0.5cm}
    \begin{center}
        \begin{tabular}{p{6cm}p{6cm}}
            \toprule
            \textbf{\large کدنویس} & \textbf{\large مهندس نرم‌افزار} \\
            \midrule
            \rule{0pt}{1em}کد می‌نویسد & مسئله را حل می‌کند \\
            \hline
            \rule{0pt}{1em}تنها کار می‌کند & در تیم کار می‌کند \\
            \hline
            \rule{0pt}{1em}به امروز فکر می‌کند & به ۵ سال بعد فکر می‌کند \\
            \hline
            \rule{0pt}{1em}تست را ضروری نمی‌داند & تست را بخشی از کار می‌داند \\
            \hline
            \rule{0pt}{1em}مشتری را نمی‌شناسد & مشتری را می‌فهمد \\
            \bottomrule
        \end{tabular}
    \end{center}
}

\slide{
    \begin{center}
        \vspace{3cm}
        {\Huge\bfseries مهندسی نرم‌افزار}
        \vspace{1.5cm}
        {\LARGE\myaccent{= ساختن چیزی که کار می‌کند،}}
        \vspace{0.5cm}
        {\LARGE\myaccent{در زمان مشخص، با تیم، و قابل نگهداری.}}
    \end{center}
}

\slide{
    \slidetitle{اهداف درس}
    \vspace{0.5cm}
    {\large در پایان این درس می‌توانید:}
    \vspace{0.5cm}
    \begin{enumerate}[itemsep=0.5em]
        \item نیازمندی‌ها را از مشتری استخراج کنید.
        \item معماری و طراحی یک سیستم را انجام دهید.
        \item در یک تیم چابک کار کنید.
        \item کد را تست، بازبینی و CI/CD تحویل دهید.
        \item کیفیت، ریسک و بدهی فنی را ارزیابی کنید.
    \end{enumerate}
}

\slide{
    \slidetitle{نقشهٔ ۱۴ هفته}
    \vspace{0.5cm}
    \begin{center}
        \begin{tabular}{cl}
            \toprule
            \textbf{هفته} & \textbf{موضوع} \\
            \midrule
            \rule{0pt}{1em}۱-۲ & نیازمندی‌ها \\
            \hline
            \rule{0pt}{1em}۳-۴ & مدل‌سازی و معماری \\
            \hline
            \rule{0pt}{1em}۵ & Sprint 1 Review \\
            \hline
            \rule{0pt}{1em}۶-۷ & طراحی شیءگرا و الگوها \\
            \hline
            \rule{0pt}{1em}۸ & میان‌ترم \\
            \hline
            \rule{0pt}{1em}۹-۱۰ & تست و CI/CD \\
            \hline
            \rule{0pt}{1em}۱۱ & Sprint 2 Review \\
            \hline
            \rule{0pt}{1em}۱۲-۱۳ & کیفیت، ریسک، DevOps \\
            \hline
            \rule{0pt}{1em}۱۴ & Demo Day \\
            \bottomrule
        \end{tabular}
    \end{center}
}

\slide{
    \slidetitle{پروژهٔ من: هوم‌استور}
    \vspace{0.5cm}
    \begin{center}
        {\Large\bfseries یک فروشگاه آنلاین لوازم خانگی}
    \end{center}
    \vspace{1cm}
    \begin{itemize}[itemsep=0.5em]
        \item من در طول ترم این پروژه را زنده می‌سازم.
        \item هدف: نشان دادن نمونهٔ عملی از هر مفهوم درس.
        \item شما هر هفته پیشرفت آن را در GitHub می‌بینید.
        \item لازم نیست کپی کنید؛ فقط اصول را یاد بگیرید.
    \end{itemize}
}

\slide{
    \slidetitle{تسک هفتهٔ آینده}
    \vspace{1cm}
    \begin{center}
        {\Large\bfseries\myaccent{تسک شما (تیم‌ها)}}
    \end{center}
    \vspace{0.5cm}
    \begin{itemize}[itemsep=0.3em]
        \item تیم ۳-۴ نفره تشکیل دهید
        \item پروژه انتخاب کنید
        \item Repository در GitHub بسازید
        \item \lr{01-vision.md} و \lr{personas.md} بنویسید
        \item ۱۰ User Story اولیه بنویسید
    \end{itemize}
}

\slide{
    \begin{center}
        \vspace{3cm}
        {\Huge\bfseries سؤالی دارید؟}
        \vspace{1.5cm}
        {\LARGE\myaccent{به امید یک ترم پربار}}
    \end{center}
}

\end{document}
\documentclass[14pt, a4paper, 
landscape]{extarticle}

% ---------- پکیج‌های پایه ----------
\usepackage{xcolor}
\usepackage{geometry}
\usepackage{graphicx}
\usepackage{booktabs}
\usepackage{array}
\usepackage{longtable}
\usepackage{titlesec}
\usepackage{enumitem}
\usepackage{tcolorbox}
\usepackage{fancyhdr}
\usepackage{tikz}

\geometry{margin=1cm}

% ---------- xepersian (آخرین پکیج) ----------
\usepackage{xepersian}
\settextfont{Vazirmatn}
\setlatintextfont{Times New Roman}
\setdigitfont{Vazirmatn}

% ---------- رنگ‌ها ----------
\definecolor{primary}{RGB}{139, 90, 43}
\definecolor{accent}{RGB}{224, 122, 63}
\definecolor{bg}{RGB}{245, 240, 232}
\definecolor{darkbg}{RGB}{92, 58, 26}

% ---------- تنظیمات صفحه ----------
\pagestyle{fancy}
\fancyhf{}
\fancyhead[R]{\small\textcolor{primary}{مهندسی نرم‌افزار — جلسهٔ ۱}}
\fancyfoot[C]{\textcolor{primary}{\thepage}}
\renewcommand{\headrulewidth}{0.5pt}
\renewcommand{\footrulewidth}{0pt}

\setlength{\headheight}{15pt}

% ---------- فاصله پاراگراف ----------
\setlength{\parindent}{0pt}
\setlength{\parskip}{0.5em}

% ---------- دستور برای اسلاید ----------
\newcommand{\slide}[1]{%
    \newpage
    \begin{tcolorbox}[
        colback=bg,
        colframe=primary,
        boxrule=1.5pt,
        arc=4mm,
        left=8mm, right=8mm, top=6mm, bottom=6mm,
        width=\textwidth,
        height=0.85\textheight
    ]
    #1
    \end{tcolorbox}
}

\newcommand{\slidetitle}[1]{%
    \begin{center}
        {\LARGE\bfseries\textcolor{primary}{#1}}
    \end{center}
    \vspace{0.5cm}
}

\newcommand{\bigtext}[1]{%
    \begin{center}
        {\Huge\bfseries #1}
    \end{center}
}

\newcommand{\myaccent}[1]{\textcolor{accent}{#1}}

% ============================================================
\begin{document}

\pagestyle{empty}

% ---------- اسلاید عنوان ----------
\begin{tikzpicture}[remember picture, overlay]
    \fill[darkbg] (current page.north west) rectangle (current page.south east);
\end{tikzpicture}

\begin{center}
    \vspace*{4cm}
    {\color{white}\Huge\bfseries مهندسی نرم‌افزار}
    
    \vspace{1cm}
    
    {\color{accent}\LARGE جلسهٔ ۱ — معرفی درس}
    
    \vspace{2cm}
    
    {\color{white}\Large ابراهیم فضلی}
    
    \vspace{0.5cm}
    
    {\color{white}\large ترم ۱۴۰۴}
\end{center}

\pagestyle{fancy}

% ---------- اسلاید ۲ ----------
\slide{
    \slidetitle{یک تصویر، یک سؤال}
    \vspace{1cm}
    \begin{center}
        {\Large\bfseries نرم‌افزار هم می‌تواند شکست بخورد.}
        
        \vspace{1.5cm}
        
        {\Huge\bfseries\myaccent{چرا؟}}
    \end{center}
}

% ---------- اسلاید ۳ ----------
\slide{
    \slidetitle{آمار تکان‌دهنده}
    \vspace{1cm}
    \begin{center}
        {\fontsize{80}{90}\selectfont\bfseries\myaccent{۷۰٪}}
        
        \vspace{1cm}
        
        {\Large پروژه‌های نرم‌افزاری شکست می‌خورند،}
        
        {\Large دیر تحویل می‌شوند یا از بودجه فراتر می‌روند.}
        
        \vspace{1.5cm}
        
        {\small منبع: Standish Group Chaos Report}
    \end{center}
}

% ---------- اسلاید ۴ ----------
\slide{
    \slidetitle{کدنویس vs مهندس نرم‌افزار}
    \vspace{0.5cm}
    \begin{center}
        \begin{tabular}{p{6cm}p{6cm}}
            \toprule
            \textbf{\large کدنویس} & \textbf{\large مهندس نرم‌افزار} \\
            \midrule
            \rule{0pt}{1em}کد می‌نویسد & مسئله را حل می‌کند \\
            \hline
            \rule{0pt}{1em}تنها کار می‌کند & در تیم کار می‌کند \\
            \hline
            \rule{0pt}{1em}به امروز فکر می‌کند & به ۵ سال بعد فکر می‌کند \\
            \hline
            \rule{0pt}{1em}تست را ضروری نمی‌داند & تست را بخشی از کار می‌داند \\
            \hline
            \rule{0pt}{1em}مشتری را نمی‌شناسد & مشتری را می‌فهمد \\
            \bottomrule
        \end{tabular}
    \end{center}
}

% ---------- اسلاید ۵ ----------
\slide{
    \begin{center}
        \vspace{3cm}
        {\Huge\bfseries مهندسی نرم‌افزار}
        
        \vspace{1.5cm}
        
        {\LARGE\myaccent{= ساختن چیزی که کار می‌کند،}}
        
        \vspace{0.5cm}
        
        {\LARGE\myaccent{در زمان مشخص، با تیم، و قابل نگهداری.}}
    \end{center}
}

% ---------- اسلاید ۶ ----------
\slide{
    \slidetitle{اهداف درس}
    \vspace{0.5cm}
    {\large در پایان این درس می‌توانید:}
    \vspace{0.5cm}
    \begin{enumerate}[itemsep=0.5em]
        \item نیازمندی‌ها را از مشتری استخراج کنید.
        \item معماری و طراحی یک سیستم را انجام دهید.
        \item در یک تیم چابک کار کنید.
        \item کد را تست، بازبینی و CI/CD تحویل دهید.
        \item کیفیت، ریسک و بدهی فنی را ارزیابی کنید.
    \end{enumerate}
}

% ---------- اسلاید ۷ ----------
\slide{
    \slidetitle{نقشهٔ ۱۴ هفته}
    \vspace{0.5cm}
    \begin{center}
        \begin{tabular}{cl}
            \toprule
            \textbf{هفته} & \textbf{موضوع} \\
            \midrule
            \rule{0pt}{1em}۱-۲ & نیازمندی‌ها \\
            \hline
            \rule{0pt}{1em}۳-۴ & مدل‌سازی و معماری \\
            \hline
            \rule{0pt}{1em}۵ & Sprint 1 Review \\
            \hline
            \rule{0pt}{1em}۶-۷ & طراحی شیءگرا و الگوها \\
            \hline
            \rule{0pt}{1em}۸ & میان‌ترم \\
            \hline
            \rule{0pt}{1em}۹-۱۰ & تست و CI/CD \\
            \hline
            \rule{0pt}{1em}۱۱ & Sprint 2 Review \\
            \hline
            \rule{0pt}{1em}۱۲-۱۳ & کیفیت، ریسک، DevOps \\
            \hline
            \rule{0pt}{1em}۱۴ & Demo Day \\
            \bottomrule
        \end{tabular}
    \end{center}
}

% ---------- اسلاید ۸ ----------
\slide{
    \slidetitle{پروژهٔ من: هوم‌استور}
    \vspace{0.5cm}
    \begin{center}
        {\Large\bfseries یک فروشگاه آنلاین لوازم خانگی}
    \end{center}
    
    \vspace{1cm}
    
    \begin{itemize}[itemsep=0.5em]
        \item من در طول ترم این پروژه را زنده می‌سازم.
        \item هدف: نشان دادن نمونهٔ عملی از هر مفهوم درس.
        \item شما هر هفته پیشرفت آن را در GitHub می‌بینید.
        \item لازم نیست کپی کنید؛ فقط اصول را یاد بگیرید.
    \end{itemize}
}

% ---------- اسلاید ۹ ----------
\slide{
    \slidetitle{تسک هفتهٔ آینده}
    \vspace{1cm}
    \begin{center}
        {\Large\bfseries\myaccent{تسک شما (تیم‌ها)}}
    \end{center}
    \vspace{0.5cm}
    \begin{itemize}[itemsep=0.3em]
        \item تیم ۳-۴ نفره تشکیل دهید
        \item پروژه انتخاب کنید
        \item Repository در GitHub بسازید
        \item \lr{01-vision.md} و \lr{personas.md} بنویسید
        \item ۱۰ User Story اولیه بنویسید
    \end{itemize}
}

% ---------- اسلاید ۱۰ ----------
\slide{
    \begin{center}
        \vspace{3cm}
        {\Huge\bfseries سؤالی دارید؟}
        
        \vspace{1.5cm}
        
        {\LARGE\myaccent{به امید یک ترم پربار}}
    \end{center}
}

\end{document}
